% BHU PYQ -- Transcribed & Corrected
% Paper: CSCSE31: Python Programming
% Exam: B.Sc. (Hons.) Semester III Examination 2025-26 (NEP)
% Category: Skill Enhancement Course (SEC)
% Department: Computer Science

\documentclass[11pt,a4paper]{article}
\usepackage[a4paper,top=2.5cm,bottom=2.5cm,left=2.8cm,right=2.8cm,headheight=14pt]{geometry}
\usepackage{amsmath,amssymb}
\usepackage[T1]{fontenc}
\usepackage[utf8]{inputenc}
\usepackage{lmodern,microtype,fancyhdr,enumitem,hyperref,lastpage,parskip,listings}

\hypersetup{
    colorlinks=true,
    urlcolor=black,
    linkcolor=black,
    pdftitle={CSCSE31: Python Programming -- BHU Sem III 2025-26 (NEP SEC)}
}

\pagestyle{fancy}
\fancyhf{}
\renewcommand{\headrulewidth}{0.4pt}
\renewcommand{\footrulewidth}{0.4pt}
\fancyhead[L]{\small \textit{Banaras Hindu University}}
\fancyhead[R]{\small \textit{CSCSE31: Python Programming (2025-26)}}
\fancyfoot[C]{\small Page \thepage\ of \pageref{LastPage}}

\newcommand{\pts}[1]{\hfill \textit{[#1]}}

\lstset{
  language=Python,
  basicstyle=\ttfamily\small,
  keywordstyle=\bfseries,
  showstringspaces=false,
  frame=single,
  frameround=tttt
}

\begin{document}

\begin{center}
  {\large\bfseries BANARAS HINDU UNIVERSITY}\\[2pt]
  {\normalsize B.Sc. (Hons.) Semester III Examination 2025-26 (NEP)}\\[6pt]
  \rule{0.8\linewidth}{0.5pt}\\[6pt]
  {\large\bfseries COMPUTER SCIENCE (SEC)}\\[4pt]
  {\large\bfseries Paper Code: CSCSE31 : Python Programming}\\[4pt]
  {\normalsize Time: 2$\frac{1}{2}$ Hours\hfill Maximum Marks: 70}
\end{center}

\rule{\linewidth}{0.4pt}

\smallskip

\noindent \textit{Instructions: All questions are compulsory. Sub-questions to be attempted are indicated in the question itself.}

\bigskip

\noindent \textbf{Question 1.} Attempt ALL of the following: \pts{2 $\times$ 5 = 10}
\begin{enumerate}[label=\alph*.]
  \item What is the difference between a list and a tuple in Python?
  \item What keyword is used to import a module in Python? Write one import example.
  \item What is the use of the \texttt{\_\_init\_\_()} method in a Python class?
  \item What does the \texttt{open()} function do in file handling? Give its full syntax.
  \item What is an f-string in Python, and how do you use it to include a variable in a string?
\end{enumerate}

\bigskip

\noindent \textbf{Question 2.} Attempt any THREE of the following: \pts{4 $\times$ 3 = 12}
\begin{enumerate}[label=\alph*.]
  \item Write a Python program to check if a number is positive, negative, or zero.
  \item Explain how the \texttt{=} operator differs from \texttt{==} in Python. Provide an example.
  \item What will be the output of the \texttt{random.randint(1, 10)} function? Explain.
  \item Explain the concept of Abstraction and provide an example in Python.
\end{enumerate}

\bigskip

\noindent \textbf{Question 3.} Attempt any THREE of the following: \pts{4 $\times$ 3 = 12}
\begin{enumerate}[label=\alph*.]
  \item Define Inheritance in Python. Write a simple example to show how a child class can use a method from a parent class.
  \item Write a Python program that accepts five numbers from the user one by one, finds their sum, and displays whether the sum is even or odd.
  \item Define a Dictionary in Python and demonstrate how to access, add, and remove elements.
  \item Explain Polymorphism in Python with a suitable example where the same method name behaves differently for different objects.
\end{enumerate}

\bigskip

\noindent \textbf{Question 4.} Attempt any THREE of the following: \pts{4 $\times$ 3 = 12}
\begin{enumerate}[label=\alph*.]
  \item What is the difference between \texttt{open('file', 'w')} and \texttt{open('file', 'a')} in Python?
  \item Differentiate between \texttt{remove()} and \texttt{pop()} in the context of Python lists with a suitable example.
  \item What is Encapsulation in OOP? Explain with an example.
  \item Write a Python program to plot a line graph showing the variation of temperature over a week using Matplotlib.
\end{enumerate}

\bigskip

\noindent \textbf{Question 5.} Attempt any TWO of the following: \pts{6 $\times$ 2 = 12}
\begin{enumerate}[label=\alph*.]
  \item What is programming? Explain the purpose of programming and describe the concepts of Variables, Input/Output (I/O), Operators, and Control Structures with suitable examples.
  \item Write a Python program that uses a recursive function to calculate and display the factorial of a number entered by the user.
  \item Explain the concept of File Handling in Python. Write a Python program to create a file, write some content, append more content, and then read and display the content from the file.
\end{enumerate}

\bigskip

\noindent \textbf{Question 6.} Attempt any TWO of the following: \pts{6 $\times$ 2 = 12}
\begin{enumerate}[label=\alph*.]
  \item Create a Python Bank Account system using Object-Oriented Programming principles. Implement methods for deposit, withdraw, and check balance.
  \item Write a Python program to create and use a user-defined module that calculates the area and perimeter of a circle.
  \item Write a Python program that accepts employee details including name, designation, and salary from the user: Format the data in rows whereby in each row, the name is left-aligned in 30 characters, the designation is left-aligned in 20 characters, and the salary is right-aligned in 10 characters with 2 decimal places. Save the formatted data to a file.
\end{enumerate}

\bigskip
\begin{center}
  \textbf{***}
\end{center}

\end{document}
